\section{Clean Events and Terminal Reward}
\label{app:clean}
\label{app:additional}
Section~\ref{sec:rl} shows that phase rotation can learn to count from a
terminal reward and extrapolate beyond the training lengths. These controls
trace that result to the input-set angle and the information carried by the
phase readout.

\subsection{Training and evaluation}
\label{app:clean-training}
The sequence policies use advantage actor-critic, a terminal reward of
$\pm1$, full-episode backpropagation, Adam,
and gradient clipping at norm~1. The default settings are width 16, batch size
256, and a learning rate of $0.003$. The fixed-length phase rotation uses an
unregularized, unnormalized state update with diagonal rotation. The length
curriculum uses normalization and the group-angle penalty from
Section~\ref{sec:stab}.

The fixed-length Transformer uses positional encoding and a learning rate
of $0.001$. It trains for $20.48$ million environment steps, compared with
$7.68$ million for the other models. The cue and monotone results confirm
that the GRU and LSTM learn these simpler memory tasks.
Table~\ref{tab:main} reports fixed-length success, median crossing steps,
and seed counts under the common budget.

\begin{table}[tbp]
\centering\small
\caption{Learning from terminal reward on clean sequences of length 10.
Final success is the mean at the end of training over five seeds.
Crossing means reaching $0.90$ greedy success on parity or $0.70$ on $\Z_3$.
The median step count includes runs censored at their training budget
($7.68$M steps for GRU, LSTM, and Born, and $20.48$M for the Transformer).
The seed count uses the common budget of $7.68$M steps.}
\label{tab:main}
\begin{tabular}{llrrrr}
\toprule
Task & Measure & GRU & LSTM & Transformer & Born \\
\midrule
Cue & Final success & 1.00 & 1.00 & 1.00 & 1.00 \\
Monotone & Final success & 1.00 & 1.00 & 1.00 & 1.00 \\
\midrule
$\Z_2$ & Final success & .50 & .50 & .883 & .901 \\
      & Median crossing steps & $>$7.68M & $>$7.68M & 15.7M & 2.6M \\
      & Seeds crossing by $7.68$M & 0/5 & 0/5 & 1/5 & 4/5 \\
$\Z_3$ & Final success & .40 & .33 & .975 & .997 \\
      & Median crossing steps & $>$7.68M & $>$7.68M & 2.4M & 0.26M \\
      & Seeds crossing by $7.68$M & 0/5 & 0/5 & 5/5 & 5/5 \\
\bottomrule
\end{tabular}
\end{table}

\FloatBarrier

The clean sequence control and the noisy feature experiments in
Appendix~\ref{app:noisy} share a trainer. Each of eight families uses two
learning rates ($0.001,0.003$) and two tuning seeds (100,101), each with 4,000 updates.
Selection averages the last three $L=14$ validation evaluations. Five independent
seeds (0--4) then run 6,000 A2C updates with batch 256, terminal reward,
and training lengths $4,6,8,10,12,14$. Each test uses 1,024 episodes per seed.

\subsection{Extrapolation across group orders}
\label{app:kboundary}
Table~\ref{tab:kscaling} compares greedy accuracy across cycle sizes after
training with terminal reward on lengths 4 through 14. Phase rotation
extrapolates best on the smaller cycles, and the LSTM is more accurate at
length 40 for $k=4,5$. The decoding margin in
Lemma~\ref{lem:margin} quantifies how tightly the ideal scores are spaced as
$k$ increases. The clean sequence control below tests exact angles
(Table~\ref{tab:newseq5}).

\begin{table}[tbp]
\centering
\caption{Zero-shot greedy accuracy at three evaluation lengths (five seeds; training curriculum
$L\in\{4,6,\dots,14\}$). Bold marks the higher accuracy of the two models.}
\label{tab:kscaling}
\small
\setlength{\tabcolsep}{3.5pt}
\begin{tabular}{cc ccc ccc}
\toprule
 & & \multicolumn{3}{c}{\textbf{Phase rotation}} & \multicolumn{3}{c}{LSTM} \\
\cmidrule(lr){3-5} \cmidrule(lr){6-8}
$k$ & chance & L10 & L40 & L80 & L10 & L40 & L80 \\
\midrule
\wsrow $2$ & 0.50 & \textbf{1.00} & \textbf{1.00} & \textbf{0.99} & 0.50 & 0.50 & 0.50 \\
\wsrow $3$ & 0.33 & \textbf{0.91} & \textbf{0.90} & \textbf{0.70} & 0.47 & 0.43 & 0.43 \\
$4$ & 0.25 & 0.80 & 0.55 & 0.54 & \textbf{0.88} & \textbf{0.85} & \textbf{0.84} \\
$5$ & 0.20 & 0.82 & 0.57 & 0.57 & \textbf{0.84} & \textbf{0.84} & \textbf{0.84} \\
\bottomrule
\end{tabular}
\end{table}


\subsection{Diagonal complex state-space baseline and readout ablations}
\label{app:lru}
\label{sec:ablation}
\label{sec:priorart}
Table~\ref{tab:mech} tests phase updates and readouts in a matched parity
comparison. Table~\ref{tab:lru} compares linear recurrent units (LRUs) on
length-10 parity over five seeds, using the same width of 16, trainer, and
budget.
Following \citet{orvieto2023lru}, the recurrence is
$x_t=\lambda\odot x_{t-1}+\gamma\odot(Bu_t)$, with
$\lambda_j=\exp(-\exp(\nu_j)+i\exp(\phi_j))$ and
$\gamma_j=\sqrt{1-|\lambda_j|^2}$. We initialize $|\lambda_j|$ uniformly
in $[0.9,0.999]$ and its angle in $[0,2\pi)$. In the unit-modulus
controls we set $|\lambda_j|=\gamma_j=1$. The input-conditioned variant in
Table~\ref{tab:lru} uses the same unit-modulus recurrence and squared-amplitude
readout with the angle set by the observation. It trains without the
stabilizers or the length curriculum.

The fixed unit-modulus rotation with a squared-amplitude readout averages
$0.66$ parity accuracy at length 10. Setting its angle from the input
raises accuracy to $0.996$ at the same width and training budget, with
$0.57$ at length 20 without stabilizers. The stabilizers target accumulating
angular error (Section~\ref{sec:stab}). Phase-advance and readout ablations
in Table~\ref{tab:mech} link the parity gains to phase carried by the
transition. The squared-amplitude readout gives the highest accuracy
among frozen-state probes on clean parity at length 80
(Appendix~\ref{app:clean-probes}).

For a single fixed linear layer, the state is affine in the input and the
difference between two quadratic scores has degree at most two. Parity
for $L>2$ is orthogonal under the uniform input distribution to every
such polynomial. A polynomial with the same sign as parity on every input
would have positive inner product with it. A small constant shift can absorb
a fixed tie rule on the finite input set. Stacking recurrent layers with a
gated linear unit changes the degree bound.
The two-layer variant reaches mean accuracy $0.72$ (best seed
$0.90$) at the training length and falls to chance by $L=20$.

\begin{table}[H]
\centering
\caption{Phase and readout ablations on parity. Entries give greedy
accuracy at $L{=}10$ under a matched trainer. The diagonal complex
state-space model (LRU) averages five seeds (Appendix~\ref{app:lru}).}
\label{tab:mech}
\small
\setlength{\tabcolsep}{4pt}
\begin{tabular}{lcc}
\toprule
Model & Parameters & Parity accuracy \\
\midrule
\multicolumn{3}{l}{\emph{Prior architectures and reductions}} \\
GRU (linear readout) & 1.1k & 0.50 \\
GRU + general quadratic readout & 1.6k & 0.50 \\
LSTM & 1.5k & 0.50 \\
uRNN (linear readout) & 0.8k & 0.50 \\
uRNN + squared-amplitude readout & 0.8k & 0.70 \\
LRU (linear readout) & 0.3k & 0.50 \\
LRU + squared-amplitude readout & 0.3k & 0.64 \\
Transformer (causal) & 20k & \textbf{0.99} \\
\midrule
\multicolumn{3}{l}{\emph{Our recurrence, different readouts}} \\
linear (softmax) readout & 0.4k & 0.50 \\
general quadratic readout & 2.4k & 0.58 \\
Squared-amplitude readout, phase $\theta\equiv 0$ & 0.3k & 0.56 \\
Squared-amplitude readout, with mixing $M$ & 0.3k & 0.72 \\
\wsrow \textbf{Squared-amplitude readout ($M{=}I$)} & \textbf{337} & \textbf{1.00} \\
\bottomrule
\end{tabular}
\end{table}


\begin{table}[H]
\centering
\caption{Diagonal complex state-space baselines on parity. Models train at $L=10$.
Entries give mean greedy accuracy over five seeds, with the best seed in
parentheses. Predicting odd only at count five gives $0.746$ at $L=10$.}
\label{tab:lru}
\small
\setlength{\tabcolsep}{2.5pt}
\begin{tabular}{lcccc}
\toprule
model & params & $L{=}10$ & $L{=}20$ & $L{=}80$ \\
\midrule
LRU + linear readout & 259 & 0.50 & --- & --- \\
LRU + general quadratic & 2.2k & 0.58 & --- & --- \\
LRU + squared-amplitude readout & 257 & 0.64 (0.76) & 0.50 & 0.50 \\
LRU, $|\lambda|{=}1$, + squared-amplitude readout & 257 & 0.66 (0.75) & 0.53 & 0.50 \\
LRU $\times2$ + GLU + squared-amplitude readout & 1.9k & 0.72 (0.90) & 0.50 & 0.50 \\
\midrule
\wsrow \textbf{+ observation-conditioned} & & & & \\
\wsrow \textbf{\quad rotation (Eq.~\ref{eq:evolve})} & \textbf{337} & \textbf{0.996} (1.00) & 0.57 & 0.52 \\
\bottomrule
\end{tabular}
\end{table}

\subsection{Generalization with terminal state labels}
\label{sec:lengthgen-arch}

Terminal state labels test extrapolation under direct supervision. Each
model observes a binary event sequence and receives a single label for the final
residue. Every model minimizes terminal-label NLL for 22,000 Adam updates at
learning rate $0.0015$, batch size 256, and lengths $4,6,\ldots,14$. The
phase-rotation model uses an angle lattice and Fourier readout initialization.

Table~\ref{tab:lengthgen-arch} compares ten seeds at the common horizon $L=400$.
The phase rotation with an angle lattice solves a majority of seeds at every
cycle size. GRU and LSTM nearly always extrapolate at $k=3$--$5$ in this
terminal-label protocol. The successful-seed counts at $L=1600$ for
$k=2,3,4,5$ are
$(9,9,6,9)$ for the phase-rotation model, $(0,10,10,9)$ for GRU,
and $(1,9,10,10)$ for LSTM. Archived predictions reproduce these counts.
\begin{table}[H]
\centering\small
\caption{Extrapolation after terminal-label supervision. Entries count
runs exceeding $0.90$ greedy accuracy at $L=400$ out of ten seeds per
configuration. Recurrent models use $3\times512$ evaluation sequences per
seed, and the Transformer uses 128.}
\label{tab:lengthgen-arch}
\begin{tabular}{lccccc}
\toprule
Architecture & Parameters & $\Z_2$ & $\Z_3$ & $\Z_4$ & $\Z_5$ \\
\midrule
Phase rotation (angle lattice, Fourier initialization) & 337--433 & 9/10 & 9/10 & 6/10 & 9/10 \\
GRU & 3.7--3.8k & 0/10 & 10/10 & 10/10 & 10/10 \\
LSTM & 5.0--5.1k & 1/10 & 10/10 & 10/10 & 10/10 \\
Transformer (no positional encoding) & 20k & 0/10 & 0/10 & 0/10 & 0/10 \\
\bottomrule
\end{tabular}
\end{table}


\subsection{Clean sequence control}
On clean inputs, a cyclic counter can carry the event count by applying exact
shifts. Table~\ref{tab:newseq5} shows that the learned counter maintains perfect
greedy accuracy across evaluation lengths for $k=5$, with the wider LSTM
close behind.
Appendix~\ref{app:adapters} describes the state-tracking adapters.

\begin{table}[H]
\centering\small
\caption{Clean sequence control for $k=5$. Greedy accuracy is the five-seed mean
with sample SD in parentheses. LR is the validation-selected learning rate.}
\label{tab:newseq5}
\setlength{\tabcolsep}{3pt}
\begin{tabular}{lccccc}
\toprule
Model & LR & $L=14$ & $L=40$ & $L=80$ & $L=160$ \\
\midrule
Learned cyclic counter & 0.001 & 1.000 (0.000) & 1.000 (0.000) & 1.000 (0.000) & 1.000 (0.000) \\
DeltaNet adapter & 0.003 & 0.388 (0.090) & 0.194 (0.010) & 0.195 (0.013) & 0.203 (0.015) \\
DeltaProduct adapter & 0.001 & 0.808 (0.228) & 0.646 (0.406) & 0.635 (0.411) & 0.571 (0.347) \\
LSTM & 0.003 & 0.476 (0.336) & 0.367 (0.354) & 0.359 (0.359) & 0.352 (0.362) \\
LSTM32 & 0.001 & 0.999 (0.001) & 0.999 (0.001) & 0.998 (0.005) & 0.994 (0.011) \\
Phase rotation, angle penalty & 0.003 & 0.791 (0.140) & 0.647 (0.346) & 0.595 (0.309) & 0.252 (0.079) \\
Phase rotation, exact angles & 0.003 & 0.946 (0.085) & 0.948 (0.077) & 0.720 (0.250) & 0.400 (0.258) \\
Phase rotation, unregularized & 0.003 & 0.712 (0.097) & 0.457 (0.233) & 0.498 (0.222) & 0.162 (0.082) \\
\bottomrule
\end{tabular}
\end{table}

\subsection{Frozen-state probes with clean events}
\label{app:clean-probes}
We freeze identical recurrence checkpoints and fit readout heads on the same
6,144 labelled examples, using identical minibatch index streams for 6,000 Adam
updates at learning rate $0.003$.
The absolute-linear head sees all real and imaginary channels. The invariant
quadratic head sees every pairwise relative-phase product and has more parameters.
Both the squared-amplitude and invariant quadratic readouts fit the clean
parity training set perfectly. Table~\ref{tab:newfrozenclean} shows the
squared-amplitude readout's clean parity advantage at $L=80$.

\begin{table}[H]
\centering\small
\caption{Frozen-state probes with clean events. Entries give mean greedy accuracy
with sample SD in parentheses over five seeds, with 2,048 test examples per seed.}
\label{tab:newfrozenclean}
\setlength{\tabcolsep}{3pt}
\begin{tabular}{lcccc}
\toprule
$k$ / head & $L=14$ & $L=40$ & $L=80$ & $L=160$ \\
\midrule
2 / Original head & 1.000 (0.000) & 1.000 (0.000) & 0.904 (0.215) & 0.484 (0.469) \\
2 / Squared-amplitude readout & 1.000 (0.000) & 1.000 (0.000) & 0.999 (0.003) & 0.586 (0.442) \\
2 / Absolute linear & 0.504 (0.132) & 0.480 (0.041) & 0.486 (0.058) & 0.508 (0.031) \\
2 / Invariant quadratic & 1.000 (0.000) & 1.000 (0.000) & 0.757 (0.408) & 0.496 (0.396) \\
3 / Original head & 0.906 (0.092) & 0.715 (0.275) & 0.630 (0.358) & 0.573 (0.352) \\
3 / Squared-amplitude readout & 0.968 (0.045) & 0.652 (0.349) & 0.596 (0.374) & 0.599 (0.365) \\
3 / Absolute linear & 1.000 (0.001) & 0.738 (0.331) & 0.588 (0.400) & 0.556 (0.412) \\
3 / Invariant quadratic & 1.000 (0.000) & 0.802 (0.257) & 0.621 (0.401) & 0.563 (0.405) \\
\bottomrule
\end{tabular}
\end{table}

\FloatBarrier
